\section{Results and Analysis}
\label{sec:results}

\subsection{Temporal Signal Characteristics}

Analysis of the DALTON dataset confirms the strong periodic structure of indoor pollutant signals that VABQ exploits. Kitchen VOC measurements exhibit sharp transient peaks reaching 800+~ppb during cooking events, with rapid decay under ventilated conditions but prolonged lingering (hours) under poor ventilation \cite{karmakar2024exploring}. Bedroom CO$_2$ accumulates overnight during sleep with AC operation, reaching median levels exceeding 1200~ppm due to closed windows and recirculated air \cite{karmakar2024indoor}. Temperature shows steady diurnal trends with relatively low variance outside of direct ventilation changes. Humidity correlates with cooking and seasonal changes, increasing noticeably during the second month of the study.

The rolling variance analysis reveals that across all sites and channels, signals remain in the low-variance regime ($\Gamma_c(t) < \theta_L = 0.15$) for 62.7\% of the total measurement time. The medium-variance regime ($\theta_L \leq \Gamma_c(t) < \theta_H = 0.55$) accounts for 24.1\% of the time, while the high-variance regime ($\Gamma_c(t) \geq \theta_H$) occupies only 13.2\%. This distribution is highly favourable for energy saving, as the majority of time is spent at 4-bit resolution.

\subsection{Bit-Width Distribution per Channel}

The adaptive bit allocation varies substantially across channels, reflecting their distinct temporal dynamics. Table~\ref{tab:bitwidth} presents the distribution.

\begin{table}[H]
\centering
\caption{Bit-width distribution and mean bit-width per sensor channel under VABQ.}
\label{tab:bitwidth}
\begin{tabular}{@{}lcccc@{}}
\toprule
\textbf{Channel} & \textbf{4-bit (\%)} & \textbf{8-bit (\%)} & \textbf{16-bit (\%)} & \textbf{Mean Bits} \\
\midrule
PM$_{2.5}$ & 54.3 & 28.9 & 16.8 & 7.25 \\
PM$_{10}$  & 56.1 & 27.4 & 16.5 & 7.09 \\
CO$_2$     & 48.2 & 31.6 & 20.2 & 7.88 \\
VOC        & 42.7 & 33.8 & 23.5 & 8.46 \\
CO         & 61.4 & 25.3 & 13.3 & 6.55 \\
NO$_2$     & 63.8 & 24.1 & 12.1 & 6.33 \\
Temperature & 78.2 & 17.1 & 4.7 & 5.06 \\
Humidity   & 71.6 & 20.8 & 7.6 & 5.52 \\
\midrule
\textbf{Average} & \textbf{59.5} & \textbf{26.1} & \textbf{14.3} & \textbf{6.77} \\
\bottomrule
\end{tabular}
\end{table}

VOC and CO$_2$ require higher average bit-widths due to their complex dynamics tied to cooking, ventilation, and occupant respiration. Temperature and humidity, being slowly varying environmental parameters, spend the majority of time at 4-bit resolution. The overall mean bit-width of 6.77 bits represents a substantial reduction from the 16-bit baseline.

\subsection{Signal Reconstruction Fidelity}

The NRMSE for each channel under different configurations quantifies the trade-off between compression and fidelity. Table~\ref{tab:nrmse} presents the results.

\begin{table}[H]
\centering
\caption{NRMSE (\%) per sensor channel across experimental configurations.}
\label{tab:nrmse}
\begin{tabular}{@{}lcccc@{}}
\toprule
\textbf{Channel} & \textbf{B1} & \textbf{B3 (8-bit)} & \textbf{B4 (Adapt.)} & \textbf{P (VABQ)} \\
\midrule
PM$_{2.5}$  & 0.00 & 1.42 & 1.68 & 1.87 \\
PM$_{10}$   & 0.00 & 1.38 & 1.61 & 1.79 \\
CO$_2$      & 0.00 & 0.96 & 1.23 & 2.14 \\
VOC         & 0.00 & 1.71 & 2.05 & 2.38 \\
CO          & 0.00 & 1.15 & 1.34 & 1.53 \\
NO$_2$      & 0.00 & 1.22 & 1.47 & 1.61 \\
Temperature & 0.00 & 0.34 & 0.41 & 0.52 \\
Humidity    & 0.00 & 0.58 & 0.72 & 0.89 \\
\midrule
\textbf{Average} & \textbf{0.00} & \textbf{1.10} & \textbf{1.31} & \textbf{1.59} \\
\bottomrule
\end{tabular}
\end{table}

All NRMSE values for the VABQ configuration remain well below the 2.5\% target, with the maximum per-channel NRMSE of 2.38\% occurring for VOC---the most dynamically complex channel. The marginal increase in NRMSE relative to static 8-bit quantisation (1.10\% $\rightarrow$ 1.59\%) is offset by the substantially greater energy savings from adaptive 4-bit allocation during stable periods.

\subsection{Classification Performance}

The eight-class activity classification results demonstrate that VABQ preserves high accuracy despite aggressive sensor-level quantisation. Table~\ref{tab:classification} presents the results.

\begin{table}[H]
\centering
\caption{Classification performance across experimental configurations.}
\label{tab:classification}
\begin{tabular}{@{}lccc@{}}
\toprule
\textbf{Configuration} & \textbf{Accuracy (\%)} & \textbf{Macro F1} & \textbf{Weighted F1} \\
\midrule
B1: Full Precision     & 97.7 & 0.968 & 0.977 \\
B2: Static PTQ         & 97.2 & 0.961 & 0.972 \\
B3: Static 8-bit       & 96.8 & 0.957 & 0.968 \\
B4: Adaptive Threshold  & 96.4 & 0.951 & 0.964 \\
\textbf{P: VABQ}       & \textbf{96.1} & \textbf{0.948} & \textbf{0.961} \\
\bottomrule
\end{tabular}
\end{table}

The VABQ configuration achieves 96.1\% overall accuracy---a 1.6 percentage point degradation from the full-precision baseline---while the macro F1-score decreases from 0.968 to 0.948. This degradation is concentrated in the ``ventilation change'' and ``electronic device use'' classes, which produce subtle, low-amplitude signal perturbations that are more susceptible to 4-bit quantisation noise. The cooking and cleaning classes, which generate strong pollutant signatures, maintain F1-scores above 0.97 even under VABQ.

\subsection{Energy Consumption Analysis}

The normalised energy consumption for each system component provides the central result of this study. Table~\ref{tab:energy} presents the breakdown.

\begin{table}[H]
\centering
\caption{Normalised energy consumption and total system saving across configurations.}
\label{tab:energy}
\begin{tabular}{@{}lccccc@{}}
\toprule
\textbf{Config.} & \textbf{Sensor} & \textbf{TX} & \textbf{Inference} & \textbf{Total} & \textbf{Saving} \\
\midrule
B1: Full Prec.   & 1.000 & 1.000 & 1.000 & 1.000 & --- \\
B2: Static PTQ   & 1.000 & 1.000 & 0.252 & 0.784 & 21.6\% \\
B3: Static 8-bit & 0.391 & 0.500 & 1.000 & 0.630 & 37.0\% \\
B4: Adapt. Thresh. & 0.547 & 0.623 & 1.000 & 0.723 & 27.7\% \\
\textbf{P: VABQ} & \textbf{0.387} & \textbf{0.423} & \textbf{0.252} & \textbf{0.354} & \textbf{64.6\%} \\
\bottomrule
\end{tabular}
\end{table}

VABQ achieves a 61.3\% reduction in sensor-level ADC energy (from 1.000 to 0.387), primarily through 4-bit allocation during the 59.5\% of time when signals are in the low-variance regime. The transmission energy reduction of 57.7\% follows directly from the reduced data volume: the mean bit-width of 6.77 bits versus the 16-bit baseline yields a proportional data volume reduction. The inference energy reduction of 74.8\% results from INT8 PTQ of the 1D-CNN \cite{smoothquant2024, palakodety2025tpu}. The compound effect produces a total system energy saving of 64.6\%.

The exponential relationship between ADC bit-width and energy is the key driver of sensor-level savings. Under the Walden FoM model (Equation~\ref{eq:walden}), switching from 16-bit to 4-bit reduces per-sample ADC energy by a factor of $2^{16}/2^4 = 4096$, though this theoretical maximum is moderated by fixed overhead circuitry to an effective factor of approximately 15$\times$ in practice \cite{andrulis2024adc, reconfigADC2025}.

\subsection{Comparison with Existing Approaches}

To contextualise the VABQ results, a comparison is drawn with published methods in Table~\ref{tab:comparison}.

\begin{table}[H]
\centering
\caption{Comparison of VABQ with existing approaches in the literature.}
\label{tab:comparison}
\begin{tabular}{@{}p{3cm}cccc@{}}
\toprule
\textbf{Method} & \textbf{Year} & \textbf{Energy Saving} & \textbf{Acc. Impact} & \textbf{Dual} \\
\midrule
Adaptive Thresh. \cite{atitallah2023autonomous} & 2023 & 17.3\% reduct. & $\pm$11.7\% & No \\
SAX-LZW \cite{aljanabi2022sax} & 2022 & $\sim$40\% data & $\sim$5\% loss & No \\
Bounded-Error \cite{azar2020bounded} & 2020 & $\sim$35\% data & Bounded & No \\
SmoothQuant \cite{smoothquant2024} & 2024 & $\sim$75\% infer. & $<$1\% loss & No \\
Q-Learning Tx \cite{qlearning2025airquality} & 2025 & 40--50\% Tx & Maintained & No \\
TinyML Ozone \cite{akhyar2025tinyml} & 2025 & Low-power & $R^2$=0.95 & No \\
EdgeMLOps \cite{edgemlops2025} & 2025 & INT8 speedup & Maintained & No \\
\textbf{VABQ (Ours)} & \textbf{2026} & \textbf{64.6\% total} & \textbf{1.6\% loss} & \textbf{Yes} \\
\bottomrule
\end{tabular}
\end{table}

The distinguishing feature of VABQ is its dual-layer operation: no prior work simultaneously optimises sensor-level ADC resolution and inference-level model precision within a unified framework. The 64.6\% total system energy saving substantially exceeds the savings achievable by either sensor-level or inference-level optimisation alone (37.0\% and 21.6\%, respectively), demonstrating the compounding benefit of joint optimisation.

\subsection{Sensitivity Analysis}

\subsubsection{Effect of Window Size $\tau$}

The rolling window duration $\tau$ controls the responsiveness of the variance estimate. Shorter windows (e.g., $\tau = 2$~min) respond quickly to transients but produce noisy variance estimates, leading to frequent bit-width switching and higher computational overhead. Longer windows (e.g., $\tau = 30$~min) provide stable estimates but may miss short-duration events such as a brief cooking episode or a fan being switched on. The empirical optimum of $\tau = 10$ minutes, consistent with the DALTON HHI formulation \cite{karmakar2024exploring}, balances responsiveness with stability, achieving the lowest combined distortion-energy metric.

\subsubsection{Effect of Weighting Parameter $\alpha$}

Varying $\alpha$ from 0 (entropy-only) to 1 (variance-only) shows that $\alpha = 0.65$ minimises the Lagrangian cost. Pure variance ($\alpha = 1$) occasionally misclassifies noisy-but-stationary signals (e.g., sensor drift) as dynamic, over-allocating bit-width. Pure entropy ($\alpha = 0$) under-responds to magnitude changes in signals that occupy few quantisation levels (e.g., temperature with slow drift). The combination captures both amplitude variation and information content, providing robust bit-width decisions across all eight channels.

\subsubsection{Effect of Threshold Values}

Varying $\theta_L$ from 0.05 to 0.30 and $\theta_H$ from 0.35 to 0.75 reveals a Pareto frontier between energy saving and reconstruction fidelity. The selected values ($\theta_L = 0.15$, $\theta_H = 0.55$) lie on the knee of this frontier, where further energy savings would incur disproportionate accuracy loss. Aggressive thresholds ($\theta_L = 0.25$, $\theta_H = 0.70$) increase the 4-bit allocation to 72.4\% and energy savings to 71.2\%, but NRMSE rises to 3.8\% and classification accuracy drops to 93.7\%.
